\ExerciseSection

\begin{enumerate}
\def\labelenumi{\arabic{enumi})}
\tightlist
\item
  Let \(f(z) = z^n + a_1z^{n-1} + \cdots + a_{n-1}z + a_n\) be a polynomial of degree \(n\) with complex coefficients and put \(f(z) = \prod_{i=1}^{n}(z - z_i)\) . Let \(r_0 = \sup_i |z_i|\) .
\end{enumerate}

\begin{enumerate}
\def\labelenumi{\alph{enumi})}
\tightlist
\item
  Show that if the real number \(r > 0\) is such that
\end{enumerate}

\[
r ^ {n} \leqslant | a _ {1} | r ^ {n - 1} + | a _ {2} | r ^ {n - 2} + \dots + | a _ {n - 1} | r + | a _ {n} |,
\]

then \(r_0 \leqslant r\) , and deduce that

\[
r _ {0} \leqslant \sup \left(\mathrm{I}, \sum_ {k = 1} ^ {n} | a _ {k} |\right).
\]

\begin{enumerate}
\def\labelenumi{\alph{enumi})}
\setcounter{enumi}{1}
\item
  Let \((\lambda_i)_{1 \leqslant i \leqslant n}\) be a finite sequence of \(n\) real numbers \(>0\) such that \(\sum_{i=1}^{n} \lambda_i^{-1} = 1\) . Show that \(r_0 \leqslant \sup_k (\lambda_k |a_k|)^{1/k}\) {[}use \(a\) {]}.
\item
  Deduce from \(a\) ) that, if the coefficients \(a_i\) are all non-zero, we have
\end{enumerate}

\[
r _ {0} \leqslant \sup \left(2 | a _ {1} |, 2 \left| \frac {a _ {2}}{a _ {1}} \right|, \dots , 2 \left| \frac {a _ {n - 1}}{a _ {n - 2}} \right|, \left| \frac {a _ {n}}{a _ {n - 1}} \right|\right).
\]

\begin{enumerate}
\def\labelenumi{\alph{enumi})}
\setcounter{enumi}{3}
\tightlist
\item
  Deduce from a) that
\end{enumerate}

\[
r _ {0} \leqslant | a _ {1} - 1 | + | a _ {2} - a _ {1} | + \dots + | a _ {n - 1} - a _ {n} | + | a _ {n} |
\]

{[}consider the polynomial \((z - 1)f(z)]\) . Hence show that, if the \(a_i\) are all real and \(>0\) , we have

\[
r _ {0} \leqslant \sup \left(a _ {1}, \frac {a _ {2}}{a _ {1}}, \dots , \frac {a _ {n - 1}}{a _ {n - 2}}, \frac {a _ {n}}{a _ {n - 1}}\right).
\]

\begin{enumerate}
\def\labelenumi{\arabic{enumi})}
\setcounter{enumi}{1}
\tightlist
\item
  An algebraic number is a complex number which is algebraic over the field \(\mathbf{Q}\) of rational numbers. Show that the ordered field \(\mathbf{B}\) of all real algebraic numbers is a maximal ordered field, but that it is not complete in the topology induced by that of \(\mathbf{R}\) {[}this topology is the same as \(\mathcal{C}_0(\mathbf{B})\) ; cf.~Chapter IV, § 3, Exercise 2 b){]}.
\end{enumerate}

¶ 3) a) Let K be a maximal ordered field, endowed with the topology \(\mathfrak{G}_{0}(\mathbf{K})\) , which is compatible with its field structure (Chapter IV, § 3, Exercise 2). Let

\[
f (\mathbf {X}) = a _ {0} \mathbf {X} ^ {n} + a _ {1} \mathbf {X} ^ {n - 1} + \dots + a _ {n}
\]

be a polynomial in K{[}X{]} which has a simple root \(\alpha\) in K. Show that, for every element \(\varepsilon > 0\) in K, there exists \(\eta > 0\) in K such that every polynomial \(g(\mathbf{X}) = b_{0}\mathbf{X}^{n} + b_{1}\mathbf{X}^{n-1} + \cdots + b_{n}\) , for which \(|a_{i} - b_{i}| \leqslant \eta\) for all i, has exactly one simple root \(\beta\) such that \(|\beta - \alpha| \leqslant \varepsilon\) . (Expand f and g in powers of X - \(\alpha\) .) Give an upper bound for \(\eta\) in terms of the \(a_{i}, \alpha\) and \(\varepsilon\) .

\begin{enumerate}
\def\labelenumi{\alph{enumi})}
\setcounter{enumi}{1}
\item
  Deduce that, if K is a maximal ordered field, its completion \(\hat{K}\) with respect to \(\mathcal{G}_{0}(K)\) (which is a field with a natural order structure: Chapter IV, § 4, Exercise 2) is a maximal ordered field.
\item
  Let \(K_{0}\) be an ordered field and let \(S = K_{0}((X))\) be the field of formal power series in one indeterminate over \(K_{0}\) ; linearly order S by taking the elements \(\geqslant 0\) in S to be o and the formal power series whose term of lowest degree has coefficient \textgreater0. Show that S, endowed with the topology \(\mathcal{G}_{0}(S)\) , is complete, but that S is not a maximal ordered field (show that the polynomials \(Y^{p} - X\) of S{[}Y{]} are irreducible for each integer p \textgreater{} I).
\end{enumerate}

\(d\) ) In the example of \(c\) ), take \(\mathbf{K}_0 = \mathbb{R}\) . Let \(\mathbf{K}\) be a maximal ordered algebraic extension of \(\mathbf{S}\) ; show that \(\mathbf{K}\) is not complete in the topology \(\mathfrak{C}_0(\mathbf{K})\) . (Embed \(\mathbf{S}\) in the field \(\mathbf{E}\) of formal power series with well-ordered exponents, ordered in the same manner as \(\mathbf{S}\) , and embed \(\mathbf{K}\) in a maximal ordered extension \(\Omega\) of \(\mathbf{E}\) . Observe that \(\mathbf{E}\) is complete, and give an example of a Cauchy sequence in \(\mathbf{K}\) which converges to an element \(f \in \mathbf{E}\) which is not algebraic over \(\mathbf{S}\) ; choose \(f\) so that there is an infinite number of S-isomorphisms of \(\mathbf{E}\) into an algebraic closure of \(\Omega\) , such that the images of \(f\) under these isomorphisms are all distinct.{]}

\begin{enumerate}
\def\labelenumi{\arabic{enumi})}
\setcounter{enumi}{3}
\tightlist
\item
  Show that every continuous homomorphism (necessarily injective) \(f\) of the topological field \(\mathbf{C}\) onto a subfield of \(\mathbf{C}\) is either the identity automorphism or the automorphism \(z \to \overline{z}\) . {[}Note that we must have \(f(x) = x\) for all \(x \in \mathbf{Q}\) , and deduce that \(f(x) = x\) for all \(x \in \mathbf{R}\) .{]} Show that there exist an infinite number of discontinuous isomorphisms of \(\mathbf{C}\) onto subfields of \(\mathbf{C}\) other than \(\mathbf{C}\) itself.
\end{enumerate}

¶ 5) a) Let K be a non-commutative division subring of the division ring of quaternions H; show that the centre Z of K is a subfield of the centre R of H. (Observe that every element of K commutes with every element of the field L generated by Z ∪ R; if Z ⊕ R, the field L would be a maximal subfield of H, and K would be contained in L, contrary to hypothesis.)

\begin{enumerate}
\def\labelenumi{\alph{enumi})}
\setcounter{enumi}{1}
\tightlist
\item
  Deduce that every isomorphism \(f\) (not necessarily continuous) of \(\mathbf{H}\) onto a division subring of \(\mathbf{H}\) is an inner automorphism \(x \to axa^{-1}\) of \(\mathbf{H}\) . {[}The restriction of \(f\) to \(\mathbf{R}\) is an isomorphism of \(\mathbf{R}\) onto a subfield of \(\dot{\mathbf{R}}\) , by virtue of \(a\) ); use Exercise 3 of Chapter IV, § 3.{]}
\end{enumerate}

